% SRC: nc04-object-pass1/REPORT.md (all numbers, camera table, object table)
\section{Live results: Nested Cinema 04}

Run \texttt{nc04-object-pass1} was the sidecar's first contact with the real scene, on an
RTX 6000 Ada, wall-time \SI{352}{\second}. The COLMAP model holds 736 registered images
across five OPENCV cameras and a 2.0M-Gaussian splat; 79 frames were selected by stratified
sampling. Camera~5, a 4K walkthrough of 120 registered frames, contributes nothing because
the package ships only the \texttt{.MOV} and no frames were extracted; video decode is out of
scope for pass~1, so this is recorded as an explicit gap. Rectify dominated wall-time
(\SI{227}{\second} of \SI{352}{\second}) through HEIC decode and full-resolution undistort of the
3200-class frames.

Six object candidates were carved (Table~\ref{tab:nc04}). Four are well-supported: the table
(most-supported, 21 views), the vintage radio, and two speakers that \texttt{carve\_components}
correctly separated as distinct 3D instances from one \texttt{speaker} mask set. Two are
problematic: \texttt{armchair chair} rests on only two views (a confidence of 1.00 is unreliable
at that count, and the label is phrase-merged), and \texttt{radio\#1} is a smaller second component
that may be a genuine second unit or a fragment. Seed crops confirm real installation objects: vintage
radios, cassette decks, a record-player case, a table. Figure~\ref{fig:seedcrops} shows
three of the evidence-scored seed crops.

\begin{table}[t]
\centering\small
\caption{nc04 object candidates. Four well-carved, two marginal.
% SRC: nc04-object-pass1/REPORT.md object table (lines 22-30)
}
\label{tab:nc04}
\setlength{\tabcolsep}{4pt}
\begin{tabular}{@{}lrrcl@{}}
\toprule
Label & Gauss. & conf & views & verdict \\
\midrule
table          & 38{,}383 & 0.80 & 21 & good (best) \\
radio \#0      & 17{,}276 & 0.87 & 16 & good \\
armchair chair & 19{,}055 & 1.00 & 2  & weak \\
speaker \#0    & 6{,}754  & 0.80 & 15 & good \\
speaker \#1    & 3{,}202  & 0.95 & 15 & good (2nd inst.) \\
radio \#1      & 6{,}286  & 0.87 & 16 & fair (fragment?) \\
\bottomrule
\end{tabular}
\end{table}

\begin{figure}[t]
\centering
\begin{subfigure}{0.31\columnwidth}
\includegraphics[width=\linewidth]{nc04-object-pass1/seed/table_00_crop.png}
\caption*{\tiny table}
\end{subfigure}\hfill
\begin{subfigure}{0.31\columnwidth}
\includegraphics[width=\linewidth]{nc04-object-pass1/seed/radio_00_crop.png}
\caption*{\tiny radio}
\end{subfigure}\hfill
\begin{subfigure}{0.31\columnwidth}
\includegraphics[width=\linewidth]{nc04-object-pass1/seed/speaker_00_crop.png}
\caption*{\tiny speaker}
\end{subfigure}
\caption{Evidence-scored seed crops from the live installation.
% SRC: runs/nc04-object-pass1/seed/*.png
}
\label{fig:seedcrops}
\end{figure}

\begin{figure}[t]
\centering
% SRC: runs/nc04-object-pass1/summary.json via scripts/make_closeout_charts.py
\includegraphics[width=\columnwidth]{nc04-object-pass1/charts/nc04_objects.pdf}
\caption{Carved candidates: Gaussian count (bars), supporting views (markers) and carve
confidence. The two below the 15-view line (armchair at 2 views, radio\#1 fragment) are the
weak pair.}
\label{fig:nc04objects}
\end{figure}

\subsection{Carve metrics on the golden set}
On the gallery golden set the full object path ran over 121 frames and scored each carved object
against 33 ground-truth object clouds. Every eval metric computes --- the primary validation goal.
Carve localises the right region but does not yet separate instances on an adversarial scene: the
gallery is wall-to-wall same-class paintings, so per-label mask unions cover most of each frame and
carve votes a large fraction of the splat foreground. Three ``objects'' of 1.5--2.4M Gaussians each
overlap ground-truth \texttt{object\_035} with high reproj-IoU (0.84--0.89) but only 3D-F1 $\approx$
0.35; smaller 3D components recover individual objects better (best 3D-F1 0.519) but many are
fragments (F1 $<$ 0.15). This contrasts with nested-cinema-04, where discrete, spatially separated
objects carved cleanly. The method is sound where objects are separated and degrades where a class
tiles the scene; per-instance masks rather than per-label unions are the primary fix.

\subsection{Per-instance carve: pass 1 to pass 2}
% SRC: runs/nc04-object-pass2/summary.json (frames, objects, rejected, timings)
Pass~2 selected all 736 registered frames (including Camera~5's walkthrough, now extracted) and
switched carve to per-instance masks. Object support rose sharply: \texttt{radio\#0} went from
17{,}276 Gaussians over 16 views to 118{,}070 over 331, \texttt{table\#0} to 39{,}155 over 250,
\texttt{speaker\#0} to 52{,}964 over 242. Nine candidates passed and three were rejected on the
min-view floor (\texttt{armchair}, \texttt{gramophone}, \texttt{mirror}, all
\texttt{insufficient\_views}) rather than kept as weak detections --- the pass-1 ``weak pair''
problem is now handled by a rejected list, not a caveat. The cost is wall-time: full-frame rectify
and identify push the pass to \SI{4662}{\second}, dominated by rectify (\SI{2297}{\second}) and
identify (\SI{1891}{\second}).

\subsection{Generation bake-off}\label{sec:objgen}
% SRC: runs/nc04-generate-pass1/*/candidates.json (Lane A); models-lane2/STATUS.md §Lane-B
Both lanes produced assets whose reserved VRAM was measured under Linux/Ada below the 32\,GB target
(Figure~\ref{fig:bakeoff}, \S\ref{sec:licence}); the Windows 11 / RTX 5090 acceptance run is still required.
\textbf{Lane~A} (TRELLIS.2, single matted crop at $1024^3$) generated 12 of 12 candidates ---
radio, table and speaker $\times$ four seeds --- with no failures. Reserved-VRAM peaks were
\SI{9.1}{\giga\byte} (radio), \SI{5.3}{\giga\byte} (table) and \SI{7.7}{\giga\byte} (speaker),
and per-seed runtimes ran \SIrange{37}{120}{\second}. \textbf{Lane~B} (ReconViaGen, textured GLB)
produced a single-view radio at \SI{15.0}{\giga\byte} and a multi-view radio (top-8 IoU-ranked
views) at \SI{16.6}{\giga\byte}; both textured GLBs are $\approx$17\,MB. Every variant stayed under 32\,GB reserved on the 48\,GB Ada
dev card, Lane~A with roughly $3.5\times$ headroom --- Linux/Ada evidence, not a Blackwell/5090 result.

A first \texttt{composed\_scene.glb} exists for the radio, placing the Lane-A and Lane-B assets
into the scene at solved transforms under an explicit basis contract (Figure~\ref{fig:composed}). Each object carries
provenance classes in \texttt{composed\_scene\_objects.json}: \texttt{geometry\_class=generated},
\texttt{material\_class=generated}, \texttt{placement\_class=carve-init}, and
\texttt{derivative\_class=composed-derivative}, so the placement is marked as carve-centroid
initialised, not yet silhouette-refined.

\begin{figure*}[t]
\centering
% SRC: runs/nc04-generate-pass1/radio/composed_proof/*.png (sidecar compose stage, real nc04 cameras)
\includegraphics[width=0.98\textwidth]{renders/composed_proof.png}
\caption{Composed-scene proof: the generated radio placed back into the real Nested Cinema 04
room and rendered from two real COLMAP cameras. Green is the Lane-A (TRELLIS.2) silhouette, blue the
Lane-B (ReconViaGen) silhouette, composited over the real rectified frame at the solved placement.
Placement is evidence-based but carve-centroid initialised (\texttt{placement\_class=carve-init}, yaw
unrefined), so the overlay is roughly rather than precisely registered --- see the carve-contamination
note in \S\ref{sec:eval-ranking}.}
\label{fig:composed}
\end{figure*}

\begin{figure}[h]
\centering
% SRC: runs/nc04-generate-pass1/*/candidates.json + models-lane2/STATUS.md via make_closeout_charts.py
\includegraphics[width=\columnwidth]{nc04-generate-pass1/charts/bakeoff_vram.pdf}
\caption{Bake-off resource envelope: reserved-VRAM peak per variant against the 32\,GB ceiling,
with runtime. Lane-A peaks are the max across four seeds; Lane-B are the textured single/multi-view
runs.}
\label{fig:bakeoff}
\end{figure}

\begin{figure}[h]
\centering
% SRC: Blender turntables of runs/nc04-generate-pass1/{radio,table,speaker} best seeds (scripts/render_turntable.py)
\includegraphics[width=\columnwidth]{renders/laneA_turntables.png}
\caption{Lane-A TRELLIS.2 objects, four turntable views each (radio, table, speaker; best seed by
mesh detail). Real generated 3D assets with baked texture; carving and single-crop generation leave
visible holes on unobserved backs.}
\label{fig:turntables}
\end{figure}

\begin{figure*}[t]
\centering
% SRC: Blender turntables, runs/nc04-generate-pass1/radio (Lane A) vs runs/nc04-laneb-pass1/radio_00_multiview (Lane B)
\includegraphics[width=0.98\textwidth]{renders/radio_bakeoff_AvB.png}
\caption{The radio reconstructed by Lane~A (TRELLIS.2 single crop, top) and Lane~B (ReconViaGen
multi-view, bottom), same four viewpoints. Observations only: Lane~A yields a closed, continuous
housing with soft surface detail; Lane~B yields separated front-panel components with visible
floaters and thin shells. No quality ranking is implied --- numeric silhouette scoring is deferred
(\S\ref{sec:eval-ranking}).}
\label{fig:bakeoffvisual}
\end{figure*}

Quality ranking across lanes is deferred: the nc04 numeric silhouette scores are not yet
reliable (the scale bracket was only just fixed via the manifest \texttt{scene\_scale}, and refined
scoring is running), so visual contact sheets are the current evidence and no numeric lane winner is
claimed. The ranking chart lands when refined scoring completes (\S\ref{sec:eval-ranking}).
